\section{Du flux d'ordres au prix gaussien : le pont volume}
\label{sec:bridge}

\subsection{Le prix comme flux d'ordres accumulé}

Sur $[t,t+\Delta]$, soit $M$ le nombre de transactions exécutées, la transaction
$k$ portant le volume signé $\epsilon_k v_k$ avec $\epsilon_k\in\{-1,+1\}$ le
sens (à l'initiative de l'acheteur ou du vendeur) et $v_k>0$ la taille. Notons
$V=\sum_k v_k$ le volume total et $Q=\sum_k \epsilon_k v_k$ le flux d'ordres
signé net.

\begin{assumption}[Impact linéaire]\label{as:impact}
Le log-prix répond linéairement au flux d'ordres signé net,
$\Delta X_{t,\Delta} = X_{t+\Delta}-X_t = \lambda\,Q$, avec $\lambda>0$ constant
sur l'intervalle.
\end{assumption}

C'est l'équilibre de \citet{kyle1985} : le teneur de marché, incapable de
distinguer le flux informé du flux non informé, fixe une règle de prix linéaire
et $\lambda$ est l'inverse de la profondeur. C'est la définition du \og prix
comme fonction du volume échangé \fg{} dont a besoin le reste de la section, et
c'est le maillon le plus faible de la chaîne ; la
section~\ref{sec:bridge-caveats} dit ce qui se passe quand il cède.

\begin{assumption}[Signes échangeables et variance proportionnelle au volume]
\label{as:signs}
Conditionnellement à $V$ et à l'information disponible en $t$, la croyance du
marché sur $Q$ n'est contrainte que par ses deux premiers moments, avec
$\mathbb{E}[Q\mid V] = m$ et $\mathrm{Var}(Q\mid V) = \varsigma^2 V$.
\end{assumption}

La forme de la variance est celle qu'on obtient lorsque les signes des
transactions sont échangeables à corrélation bornée et que les tailles ont une
moyenne stable : pour des signes indépendants,
$\mathrm{Var}(Q)=\sum_k v_k^2 \approx \bar v V$. Le contenu économique est que
\emph{le volume, et non le temps, est la mesure naturelle de ce qui peut se
produire}.

\subsection{Maximum d'entropie en temps-volume}

\begin{proposition}[Maximum d'entropie sous contraintes de moments d'ordre un et deux]
\label{prop:maxent}
Parmi les distributions absolument continues sur $\mathbb{R}^N$ de moyenne $\mu$
et de covariance $\Sigma\succ0$, la gaussienne $\mathcal N(\mu,\Sigma)$ maximise
de façon unique l'entropie différentielle, et
$h^{\max}(\Sigma)=\tfrac12\log\left((2\pi e)^N\det\Sigma\right)$.
\end{proposition}

\begin{proof}
Soit $q=\mathcal N(\mu,\Sigma)$ et soit $p$ de mêmes deux premiers moments et
d'entropie différentielle finie. Comme $\log q$ est exactement quadratique, son
espérance ne dépend que de ces moments, donc
$\mathbb{E}_p[\log q]=\mathbb{E}_q[\log q]=-h(q)$. La positivité de l'entropie
relative donne $0\le \KL{p}{q} = -h(p)+h(q)$, d'où $h(p)\le h(q)$, avec égalité
si et seulement si $p=q$ presque partout. L'argument n'exige aucune
différentiabilité de $p$ ni l'existence préalable d'un maximiseur.
\end{proof}

La voie variationnelle --- un lagrangien sur les contraintes de normalisation, de
moyenne et de covariance, dont la condition du premier ordre rend $\log p$
quadratique --- explique \emph{pourquoi} le maximiseur est gaussien, mais c'est
l'argument par la divergence qui établit l'optimalité et l'unicité.

En appliquant la proposition~\ref{prop:maxent} à
l'hypothèse~\ref{as:signs} en dimension un, puis en transportant par
l'hypothèse~\ref{as:impact} :

\begin{equation}
\label{eq:mdh}
\Delta X_{t,\Delta}\ \big|\ V \ \sim\ \mathcal{N}\!\left(\lambda m,\ \lambda^2\varsigma^2 V\right).
\end{equation}

L'équation~\eqref{eq:mdh} est le pont. Trois points méritent d'être soulignés.

D'abord, \emph{la variable de conditionnement est le volume, pas le temps}. Le
maximum d'entropie ne délivre pas une variance ; il délivre la distribution la
moins engageante étant donné une variance. L'hypothèse~\ref{as:signs} fournit
cette variance à partir d'un observable, et cet observable est le volume échangé.
L'équation~\eqref{eq:mdh} est l'hypothèse de mélange de distributions de
\citet{clark1973} et \citet{tauchen1983}, obtenue ici comme un énoncé de maximum
d'entropie plutôt que posée ; \citet{ane2000} montrent empiriquement que les
rendements mesurés en temps-transaction sont proches du gaussien alors que les
rendements en temps calendaire ne le sont pas, ce qui est la prédiction
de~\eqref{eq:mdh}.

Ensuite, \emph{la gaussianité est conditionnelle}. Inconditionnellement,
$\Delta X$ est un mélange d'échelle de lois normales sur le volume aléatoire $V$,
donc leptokurtique même si \eqref{eq:mdh} est exacte. C'est le mécanisme de
subordination de \citet{mandelbrot1967}, et cela signifie que les queues épaisses
ne constituent pas à elles seules une preuve contre le maximum d'entropie.

Enfin, l'extension multivariée est immédiate : avec $N$ actifs,
$\Delta \mathbf{X}\mid V \sim \mathcal N(\boldsymbol\mu V, \Sigma V)$, et la
proposition~\ref{prop:maxent} s'applique telle quelle.

\subsection{Des incréments gaussiens au mouvement brownien puis au MBG}

Soit $\theta_t=\int_0^t v_s\,ds$ l'horloge de volume cumulé.
L'équation~\eqref{eq:mdh} donne la loi d'un incrément ; un processus demande
davantage.

\begin{assumption}[Cohérence temporelle]\label{as:time}
Dans l'horloge de volume, les incréments sur des intervalles disjoints sont
indépendants, les trajectoires sont continues, et la covariance d'un incrément
est proportionnelle au volume écoulé.
\end{assumption}

Sous l'hypothèse~\ref{as:time} à coefficients constants, il existe une matrice
$B$ telle que $BB^{\!\top}=\Sigma$ et $X_\theta = X_0 + \mu\theta + BW_\theta$ en
temps-volume, donc en temps calendaire
\begin{equation}
\label{eq:subordination}
X_t = X_0 + \mu\,\theta_t + B\,W_{\theta_t},
\end{equation}
un mouvement brownien \emph{subordonné} à l'horloge de volume. Si l'intensité de
négociation est constante, $\theta_t=\bar v\,t$ et \eqref{eq:subordination} est
un mouvement brownien ordinaire en temps calendaire ; en exponentiant et en
appliquant Itô,
\begin{equation}
\frac{dS_{i,t}}{S_{i,t}} = \left(a_{i,t}+\tfrac12\sigma_{i,t}^2\right)dt + \sigma_{i,t}\,dW_{i,t},
\end{equation}
de sorte que la dérive du prix n'est pas celle du log-prix --- la même correction
du second ordre qui reparaîtra à la section~\ref{sec:channels} lorsqu'on
différenciera $\log\det R_t$.

Il vaut la peine d'être précis sur ce que fait l'hypothèse~\ref{as:time}. Par le
théorème de \citet{monroe1978}, \emph{toute} martingale locale continue est un
mouvement brownien changé de temps : la subordination seule n'a donc aucun
contenu empirique. Le contenu est dans l'horloge : \eqref{eq:mdh} dit que la
bonne horloge est le volume, et l'hypothèse~\ref{as:time} dit que le marché est
sans mémoire dans cette horloge. C'est le maximum d'entropie qui sélectionne le
mouvement brownien dans la classe qu'autorise le théorème de Monroe.

\subsection{La contraposée}
\label{sec:contrapositive}

La proposition~\ref{prop:maxent} se lit dans les deux sens. Notons $q_t$ la
gaussienne de mêmes moyenne et covariance conditionnelles que $p_t$. Alors

\begin{equation}
\label{eq:deficit}
\DD_t \;\equiv\; h(q_t) - h(p_t) \;=\; \KL{p_t}{q_t} \;\ge\; 0,
\end{equation}
avec égalité si et seulement si $p_t$ est gaussienne. Donc : \emph{les prix sont
non gaussiens exactement dans la mesure où leur distribution conditionnelle n'est
pas d'entropie maximale}, et $\DD_t$ mesure cette mesure. Dans le langage de la
section~\ref{sec:bridge}, un $\DD_t$ positif signifie que des contraintes
au-delà des deux premiers moments sont actives --- que quelqu'un sait quelque
chose que le moment d'ordre deux n'encode pas.

Deux sources produisent un déficit positif, et les séparer importe. Une horloge
de volume aléatoire produit un mélange d'échelle, symétrique et leptokurtique.
Une information directionnelle contraignante produit de l'asymétrie. La
section~\ref{sec:channels} scinde $\DD$ exactement selon cette ligne.

\subsection{Où le pont est faible}
\label{sec:bridge-caveats}

L'hypothèse~\ref{as:impact} est linéaire ; l'impact empirique d'un métaordre est
plus proche d'une racine carrée de la taille, ce qui rend concave l'application
du flux au prix et la distribution de rendements induite moins épaisse en queue
que le flux. L'hypothèse~\ref{as:signs} traite les signes comme échangeables ;
le flux d'ordres est en réalité fortement autocorrélé, ce qui gonfle
$\mathrm{Var}(Q)$ au-dessus de $\varsigma^2 V$ et fait de l'horloge de volume un
substitut imparfait de la variance. Aucune de ces défaillances ne change la
\emph{forme} de l'argument --- le maximum d'entropie sous budget de variance
reste gaussien --- mais toutes deux signifient que l'horloge de volume est une
approximation de l'horloge d'information plutôt que cette horloge même. Nous
utilisons le raisonnement en temps-volume pour motiver le cadre et l'estimation
en temps calendaire dans tout le travail empirique, de sorte qu'aucun résultat
ci-dessous ne dépend de l'exactitude de \eqref{eq:mdh}.

\section{Les trois canaux de l'entropie}
\label{sec:channels}

\subsection{Échelle et dépendance}

Écrivons $\Sigma_t = D_tR_tD_t$ avec
$D_t=\mathrm{diag}(\sigma_{1,t},\dots,\sigma_{N,t})$. Alors
$\det\Sigma_t = \left(\prod_i\sigma_{i,t}^2\right)\det R_t$ et
\begin{equation}
\label{eq:ceiling}
h(q_t) \;=\; \underbrace{\tfrac{N}{2}\log(2\pi e)}_{\text{constante}}
\;+\; \underbrace{\sum_i \log\sigma_{i,t}}_{H^{vol}_t}
\;+\; \underbrace{\tfrac12\log\det R_t}_{H^{dep}_t}.
\end{equation}
Par l'inégalité de Hadamard, $H^{dep}_t\le 0$, avec égalité si et seulement si
$R_t=I$ : \emph{la corrélation détruit toujours de l'entropie}, et $-H^{dep}_t$
est précisément la multi-information gaussienne (corrélation totale) du vecteur
de rendements. Une lecture utile de la même quantité est spectrale : avec
$\lambda_k$ les valeurs propres de $R_t$ et $\tilde\lambda_k=\lambda_k/N$, le
nombre de participation
$N^{\text{eff}}_t=\exp\!\left(-\sum_k\tilde\lambda_k\log\tilde\lambda_k\right)$
est le nombre effectif de modes de risque indépendants ; il vaut $N$ pour $R=I$
et tend vers $1$ lorsqu'un unique mode collectif absorbe la variance.

L'équation~\eqref{eq:ceiling} est le plafond, pas l'entropie. C'est ce
qu'implique la seule matrice de covariance, et la proposition~\ref{prop:maxent}
dit que la vérité se situe en dessous.

\subsection{Forme : le déficit d'entropie et sa composition}

En soustrayant \eqref{eq:deficit} de \eqref{eq:ceiling} on obtient la
décomposition de l'article, trois canaux et une constante :
\begin{equation}
\label{eq:three}
h(p_t) \;=\; \tfrac{N}{2}\log(2\pi e) \;+\; \underbrace{H^{vol}_t}_{\text{échelle}} \;+\; \underbrace{H^{dep}_t}_{\text{dépendance}} \;-\; \underbrace{\DD_t}_{\text{forme}}.
\end{equation}

Le canal de forme, ce sont l'asymétrie et les queues épaisses, et le
développement d'Edgeworth le dit explicitement : pour une variable standardisée
d'asymétrie $S$ et de kurtosis excédentaire $K$ au voisinage de la gaussianité,
\begin{equation}
\label{eq:edgeworth}
J \;\approx\; \frac{S^2}{12} + \frac{K^2}{48}.
\end{equation}
L'équation~\eqref{eq:edgeworth} répond à la question \og où un saut d'entropie
serait-il visible ? \fg{} : dans les moments d'ordre trois et quatre, de manière
additive, avec un canal d'asymétrie et un canal de queue. Elle est aussi la
source du problème de signe que la section~\ref{sec:dynamics} doit résoudre, car
$S$ y entre au carré --- \emph{la néguentropie ne distingue pas un krach d'une
envolée}. Comme estimateur, \eqref{eq:edgeworth} est inutilisable sur des
rendements quotidiens ; la section~\ref{sec:program} documente à quel point et
ce que nous utilisons à la place.

\subsection{L'indice structurel}

Le canal d'échelle est une nuisance pour notre propos : il domine numériquement
le niveau de \eqref{eq:three}, il bouge avec la volatilité, et un changement
d'unités le modifie. Les deux autres canaux sont sans échelle. En les combinant,
et en utilisant la règle de chaînage exacte
$\KL{p}{\prod_i q_i} = \sum_i \KL{p_i}{q_i} + I(x)$ où $I$ est la
multi-information :

\begin{definition}[Indice d'information structurelle]
\label{def:J}
\begin{equation}
\JJ_t \;\equiv\; \DD_t \;-\; H^{dep}_t \;=\; \KL{p_t}{\textstyle\prod_i q_{i,t}} \;\ge\; 0 ,
\end{equation}
la divergence de la loi jointe conditionnelle des rendements à un produit de
marginales gaussiennes de mêmes échelles.
\end{definition}

$\JJ$ est nul si et seulement si les rendements sont indépendants et gaussiens ;
il s'élève quand les actifs se couplent, quand les marginales s'épaississent ou
quand elles s'asymétrisent. Il est mesuré en nats, il est invariant au niveau de
volatilité, et il est additif sur ses deux canaux, de sorte qu'on peut toujours
demander quelle part d'un mouvement vient de la dépendance et quelle part de la
forme. C'est la variable d'état de la section~\ref{sec:dynamics}.

\section{Dynamique de l'information : un théorème, deux postulats}
\label{sec:dynamics}

\subsection{Les sauts vers le bas sont un théorème}

Deux énoncés distincts sous-tendent la moitié descendante de la dynamique
postulée. Les confondre est une erreur facile, et les versions antérieures de ce
travail la commettaient ; nous les séparons donc explicitement.

Rappelons d'abord la définition de l'\emph{information mutuelle} entre le
vecteur de rendements $X$ et un signal $Y$ :
\begin{equation}
\label{eq:mi}
I(X;Y) \;=\; \KL{p_{XY}}{p_X\otimes p_Y}
\;=\; \iint p(x,y)\,\log\frac{p(x,y)}{p(x)p(y)}\,dx\,dy ,
\end{equation}
c'est-à-dire la divergence de la loi jointe au produit des marginales : de
combien la loi jointe s'écarte de l'indépendance. Elle est positive, symétrique,
nulle si et seulement si $X\perp Y$, et --- contrairement à $h$ --- invariante par
reparamétrisation inversible. Pour un couple gaussien scalaire de corrélation
$\rho$, $I=-\tfrac12\log(1-\rho^2)$.

\begin{proposition}[Chute de l'entropie conditionnelle moyenne]
\label{prop:conditioning}
Soit $X$ le vecteur de rendements et $Y$ un signal. Alors
\begin{equation*}
\mathbb{E}_Y\!\left[h(X\mid Y)\right] \;=\; h(X) - I(X;Y) \;\le\; h(X),
\end{equation*}
avec égalité si et seulement si $X\perp Y$.
\end{proposition}

\begin{proposition}[Monotonie du plafond sous contrainte]
\label{prop:constraint}
Soient $\mathcal{C}\subseteq\mathcal{C}'$ deux ensembles de contraintes et
$\mathcal{Q}(\cdot)$ l'ensemble des lois qui les satisfont. Alors
$\mathcal{Q}(\mathcal{C}')\subseteq\mathcal{Q}(\mathcal{C})$ et donc
\begin{equation*}
\sup_{p\,\in\,\mathcal{Q}(\mathcal{C}')} h(p) \;\le\; \sup_{p\,\in\,\mathcal{Q}(\mathcal{C})} h(p).
\end{equation*}
\end{proposition}

Les deux résultats sont classiques \citep{cover2006} et ne sont pas
équivalents ; la démonstration détaillée et un exemple gaussien chiffré figurent
en annexe~\ref{app:mi}. Leurs portées diffèrent sur deux axes.

La proposition~\ref{prop:conditioning} \emph{quantifie} la baisse --- exactement
$I(X;Y)$ nats --- mais seulement \emph{en moyenne} sur les réalisations du
signal. Pour un $y$ particulier, $h(X\mid Y=y)$ peut excéder $h(X)$ : un signal
suffisamment déroutant laisse le marché plus incertain qu'avant. Il faut donc
s'attendre à une composante de saut principalement, et non exclusivement,
unilatérale --- ce que trouve la section~\ref{sec:program}.

La proposition~\ref{prop:constraint} vaut au contraire \emph{pour chaque
réalisation}, mais porte sur le \emph{plafond} de maximum d'entropie et non sur
l'entropie réelle, et elle ne quantifie rien : elle dit que le plafond ne peut
pas monter, pas de combien il descend.

Ce que les deux disent ensemble, et qui motive tout ce qui suit : la chute est
\emph{discontinue}, car le signal arrive en un point du temps, et sa taille
espérée en nats est le contenu informationnel de la nouvelle --- une amplitude
mesurée sur une échelle qui ne dépend ni de l'actif, ni de la devise, ni des
unités.

\paragraph{Une réserve d'honnêteté.} Aucune de ces deux propositions n'est
utilisée pour calculer quoi que ce soit dans la partie empirique. Nous mesurons
$\JJ$, et non $h(X)-\mathbb{E}_Y[h(X\mid Y)]$ : estimer une entropie
différentielle en dimension~20 sur une fenêtre de 252~jours est hors de portée
(annexe~\ref{app:units}). Le lien entre le théorème et les statistiques
rapportées plus loin est un raisonnement, pas une identité estimée, et le
lecteur doit le tenir pour tel.

\subsection{Deux postulats}

\begin{postulate}[Relaxation]
\label{p:relax}
En l'absence d'information nouvelle, l'indice structurel décroît vers un niveau
de long terme à taux constant :
$d\JJ_t = -\kappa\left(\JJ_t-\bar\JJ\right)dt + \sigma_\JJ\,dW_t$ avec $\kappa>0$.
\end{postulate}

L'économie derrière P\ref{p:relax} est qu'une contrainte s'évente. La nouvelle
est intégrée aux prix, le désaccord se reforme, les positions se dispersent, et
la coupe transversale recouvre la liberté que le choc lui avait retirée. Rien
dans la théorie de l'information ne délivre cela ; c'est l'hypothèse substantielle
de l'article, et $\kappa$ est l'objet à estimer.

\paragraph{Pourquoi un processus d'Ornstein--Uhlenbeck et pas autre chose ?}
Le choix n'est pas un choix de commodité : il est \emph{forcé} dès qu'on impose
trois exigences minimales, et il est le seul à les satisfaire toutes les trois.

\begin{enumerate}
\item \emph{Markov.} L'avenir de $\JJ$ ne dépend du passé qu'à travers son
niveau présent. C'est l'exigence de parcimonie : ne postuler aucune mémoire
au-delà de ce que l'état courant contient déjà.
\item \emph{Stationnaire.} $\JJ$ possède un niveau de long terme et une variance
finis. Sans cela l'indice dérive, il n'est pas une variable d'état, et les
comparaisons de régime de la section~\ref{sec:program} n'ont pas de sens.
\item \emph{Bruit gaussien.} Entre deux arrivées d'information, on ne prétend
rien savoir du bruit au-delà de sa variance --- c'est le principe de maximum
d'entropie appliqué au bruit lui-même, exactement comme il l'a été aux
rendements en section~\ref{sec:bridge}.
\end{enumerate}

Le théorème de \citet{doob1942} énonce que \textbf{le processus d'Ornstein--Uhlenbeck
est l'unique processus à la fois markovien, gaussien et stationnaire} en temps
continu. Les trois exigences ne laissent donc aucune latitude sur la partie
diffusive : elle est déterminée. Le terme de saut, lui, est ajouté précisément
pour \emph{briser} la gaussianité là où le cadre affirme qu'elle ne tient pas,
c'est-à-dire aux arrivées d'information. La structure de \eqref{eq:ou} est ainsi
un squelette imposé plus une composante délibérément non gaussienne, ce qui est
exactement la classe introduite par \citet{bns2001}.

\paragraph{Ce qui, en revanche, reste arbitraire.} Trois choses, qu'il faut
énoncer plutôt que laisser croire qu'elles découlent du théorème. Le taux
$\kappa$ est supposé \emph{constant} : rien n'interdit qu'une crise soit
absorbée plus vite, ou plus lentement, qu'une nouvelle ordinaire, et les données
employées ici ne permettent pas de trancher. Le rappel est \emph{linéaire} en
l'écart au niveau de long terme, alors qu'un rappel accéléré loin de la moyenne
serait tout aussi plausible. Enfin, $\JJ$ est positif par construction alors
qu'un Ornstein--Uhlenbeck muni d'un terme brownien peut devenir négatif : la
spécification de \citet{bns2001} évite cet écueil en supprimant le terme
brownien et en pilotant le processus par un subordinateur à sauts positifs, ce
qui garantit une loi stationnaire portée par $\mathbb{R}^+$. Empiriquement la
question ne se pose pas --- $\JJ$ reste compris entre $2{,}17$ et $21{,}13$~nats
sur le panel principal et entre $0{,}42$ et $1{,}69$ sur le panel de
réplication, donc très loin de la frontière --- mais formellement il faudrait
soit supprimer $\sigma_\JJ$, soit écrire le processus sur $\log\JJ$.

Le postulat n'est enfin pas anodin sur le plan de la mesure : un estimateur doté
de mémoire fabriquera une décroissance même en l'absence de toute relaxation
économique, aussi la section~\ref{sec:program} mesure-t-elle la décroissance
qu'un monde sans mémoire produirait déjà.

\begin{postulate}[Asymétrie de la contrainte]
\label{p:sign}
La contrainte qu'impose l'information est unilatérale dans la direction des
pertes, de sorte qu'un saut d'entropie s'accompagne d'une asymétrie négative du
facteur de marché.
\end{postulate}

\paragraph{Pourquoi ce postulat est logiquement nécessaire.} L'indice $\JJ$ est
une fonction \emph{paire} de l'asymétrie : changer le signe de l'asymétrie ne
change pas sa valeur. La raison est immédiate. Le canal de forme est
$\Dun(z)=\KL{p_z}{\varphi}$, et la loi normale $\varphi$ est symétrique ; donc
si l'on remplace $z$ par $-z$, la loi de référence est inchangée et la
divergence l'est aussi :
\begin{equation}
\label{eq:even}
\Dun(-z)=\Dun(z),\qquad\text{et par conséquent}\qquad \JJ(-x)=\JJ(x).
\end{equation}
Le développement d'Edgeworth \eqref{eq:edgeworth} le rend visible à l'\oe il :
l'asymétrie $S$ y entre au carré.

La conséquence est structurelle et non anecdotique. \textbf{Une distribution à
longue queue à gauche --- un krach --- et son image en miroir à longue queue à
droite --- une envolée --- ont exactement le même $\JJ$.} L'indice mesure la
\emph{taille} de l'information, jamais son sens. Aucun raffinement de
l'estimateur n'y changera quoi que ce soit : c'est une propriété de l'objet
mesuré. Si l'article veut dire quoi que ce soit sur la direction, il doit donc
ajouter un ingrédient extérieur. P\ref{p:sign} est cet ingrédient, et c'est le
plus petit possible : une seule affirmation, sur le signe, ajoutée à un cadre
qui par construction n'en produit aucune.

\paragraph{Sur quoi il s'appuie.} Non pas sur un mécanisme que nous
affirmerions, mais sur une régularité empirique établie : l'\emph{effet de
levier} au sens de \citet{black1976}, c'est-à-dire le fait que les rendements
négatifs sont suivis de hausses de volatilité plus fortes que les rendements
positifs de même amplitude. C'est l'un des faits stylisés les plus robustes de
la finance empirique \citep{christie1982,engle1993,nelson1991}. Traduit dans le
langage de cet article : à quantité d'information égale, la réponse du marché
est plus ample vers le bas que vers le haut, donc un saut de $\JJ$ arrive
accompagné d'une asymétrie de marché négative.

Nous nous abstenons délibérément de choisir entre les explications causales
concurrentes de cet effet --- levier financier à proprement parler, rétroaction
de la volatilité, ou retrait asymétrique de la liquidité du côté vendeur. Le
postulat ne dépend d'aucune d'entre elles : il n'affirme que l'asymétrie, pas sa
cause. Il est testable séparément (H5) et réfutable séparément, et la
section~\ref{sec:program} conclut d'ailleurs qu'il n'est pas établi.

\subsection{La dynamique qui en résulte}

La proposition~\ref{prop:conditioning} jointe à P\ref{p:relax} donne un processus
d'Ornstein--Uhlenbeck non gaussien à sauts unilatéraux, de la classe introduite
par \citet{bns2001} pour la volatilité stochastique :
\begin{equation}
\label{eq:ou}
d\JJ_t \;=\; -\kappa\left(\JJ_t-\bar\JJ\right)dt \;+\; \sigma_\JJ\,dW_t \;+\; \int_{z>0} z\,\mathcal{P}(dt,dz),
\end{equation}
Le terme intégral demande d'être lu avec soin, car sa notation est compacte.

\paragraph{Ce que $\mathcal{P}$ n'est pas.} Ce n'est ni une loi gaussienne, ni
une fonction. La lettre $\mathcal N$ est réservée dans cet article à la loi
normale, et la mesure de saut porte donc un symbole distinct.

\paragraph{Ce que $\mathcal{P}$ est.} Une \emph{mesure aléatoire de comptage},
dite mesure aléatoire de Poisson. On peut se la représenter comme un semis
aléatoire de points dans le plan $(\text{date},\ \text{taille})$ : chaque point
est un événement informationnel, son abscisse est la date à laquelle il survient
et son ordonnée la taille du saut qu'il provoque. Formellement,
$\mathcal{P}(dt,dz)$ vaut $1$ s'il existe un saut dans l'intervalle de temps
$[t,t+dt)$ dont la taille tombe dans $[z,z+dz)$, et $0$ sinon. L'intégrale
$\int_{z>0} z\,\mathcal{P}(dt,dz)$ se lit alors simplement : \og s'il y a un
saut à la date $t$, ajouter sa taille $z$ à $\JJ$ ; sinon n'ajouter rien \fg{}.
C'est une écriture condensée de $\sum_{j} z_j$, la somme des sauts survenus.

\paragraph{Ce qui la caractérise.} Deux ingrédients, réunis dans sa
\emph{mesure d'intensité} $\nu(dz)\,dt$, qui est le nombre espéré de sauts de
taille comprise entre $z$ et $z+dz$ pendant $dt$ : d'une part une
\emph{fréquence d'arrivée} $\lambda=\nu(\mathbb{R}^+)$, en nombre de sauts par
unité de temps, d'autre part une \emph{loi des tailles} $\nu(dz)/\lambda$. Sur
le panel principal, la section~\ref{sec:program} estime $\lambda\approx5$ sauts
par an, une taille moyenne de $0{,}62$~nat et un maximum de $3{,}21$~nats.

\paragraph{Ce qu'est $z$.} La taille d'un saut individuel, exprimée en nats.
Conceptuellement, c'est l'information mutuelle de la nouvelle qui l'a provoqué,
au sens de la proposition~\ref{prop:conditioning}. Opérationnellement, c'est la
variation de $\JJ$ détectée par le filtre à quatre sigmas de la
section~\ref{sec:program}. Les deux ne coïncident qu'approximativement, et le
papier ne prétend pas les identifier.

\paragraph{Pourquoi l'intégrale ne porte que sur $z>0$.} Parce que
$\JJ$ mesure un \emph{déficit} d'entropie et que, par la
proposition~\ref{prop:conditioning}, l'information ne peut qu'accroître ce
déficit en espérance. Dans ce cadre, rien ne crée d'\emph{augmentation soudaine}
d'entropie : l'entropie n'est restituée que par la relaxation lente du premier
terme, jamais par un saut. D'où des sauts unilatéraux.

Deux réserves, toutefois. La proposition~\ref{prop:conditioning} est un énoncé
\emph{en moyenne} : une réalisation particulièrement déroutante peut accroître
l'entropie, donc une spécification fidèle autoriserait une petite probabilité de
sauts négatifs plutôt que de les interdire. Et la mesure de Poisson suppose des
arrivées indépendantes, alors que l'étude d'événement de la
section~\ref{sec:program} montre un regroupement net des sauts ; un processus
auto-excitant de type Hawkes serait plus fidèle.

Notons enfin que le coefficient de diffusion est noté $\sigma_\JJ$ et non
$\eta$, cette dernière lettre étant réservée au rayon des scénarios de stress
(section~\ref{sec:stress}). Lu sur l'entropie elle-même plutôt que sur
$\JJ$, le signe s'inverse : l'entropie chute par sauts et remonte par diffusion.
Les prédictions observables sont nettes, et ce sont elles que teste la
section~\ref{sec:program} : la loi de $\Delta\JJ$ est asymétrique à droite ;
les grands mouvements sont à la hausse et les petits à la baisse ; et après un
saut positif l'indice décroît avec une demi-vie $\ln 2/\kappa$.

\subsection{Pourquoi queues épaisses, corrélation et sauts arrivent ensemble}

Le cadre produit une prédiction qu'aucun modèle de volatilité ni de corrélation
ne produit seul. Elle se déduit en quatre pas, sans vocabulaire technique.

\begin{enumerate}
\item Une nouvelle importante concerne \emph{tous} les actifs à la fois : c'est
la même information qui contraint chacun d'eux.
\item Si la nouvelle est mauvaise, tous les actifs baissent donc \emph{le même
jour}.
\item Or ce jour-là apparaît aussi, pris isolément, dans l'historique de chaque
actif : c'est pour lui une journée de perte inhabituellement grande.
\item Accumuler de telles journées déforme la distribution des rendements de
chaque actif : le côté des pertes s'allonge et s'alourdit davantage que le côté
des gains.
\end{enumerate}

Le pas 2 est ce que mesure le canal de dépendance : les actifs bougent ensemble,
la corrélation mesurée monte, $-H^{dep}$ monte. Le pas 4 est ce que mesure le
canal de forme : des pertes extrêmes plus nombreuses et plus grandes que les
gains extrêmes, c'est-à-dire un kurtosis plus élevé et une asymétrie négative,
donc $\DD$ monte. Les deux canaux montent ensemble parce qu'ils décrivent le
même jour vu de deux endroits différents --- la coupe transversale pour l'un, la
distribution de chaque actif pour l'autre.

En marché calme, les mouvements marquants sont propres à chaque actif : ils
n'alignent pas la coupe transversale et n'ont donc aucune raison de faire monter
les deux canaux en même temps.

La forme testable est un \emph{couplage dépendant du régime} : la corrélation
entre $-\Delta H^{dep}$ et $\Delta\DD$ devrait être proche de zéro en marché
calme et fortement positive sous stress. C'est H6, et c'est là que se trouve la
réponse de l'article à la question \og le lien entre queues épaisses, corrélation
et sauts d'entropie exige-t-il un postulat supplémentaire ? \fg{} : non. Étant
donné la proposition~\ref{prop:conditioning}, le couplage découle du caractère
commun de la nouvelle. Seul le \emph{signe} exige P\ref{p:sign}.

